\documentclass[../../main.tex]{subfiles}
\graphicspath{{../../image/}{../../figure/}}
\begin{document}

\section{源部分耗散时的 $L^1$ 稳定性}\label{section:ch16-部分耗散源l1稳定}

此处的目标是证明\cref{theorem:ch16-1631}的假设——从而其结论——适用于刻画松弛现象的系统。因此本节是第 5.5 节的对应物，其中 $BV$ 弱解扮演经典解的角色。

如同第 5.5 节那样，我们假设系统\eqref{align:ch16-s02-e01}配有光滑的熵--熵流对 $(\eta,q)$，其中 $\eta$ 凸且已由 $\eta(0)=0$、$\mathrm{D}\eta(0)=0$ 归一化。于是任一可容 $BV$ 解满足熵不等式\eqref{align:ch16-s05-e01}。我们进一步假设，源 $P$ 相对于 $\eta$ 是耗散半正定的，用第 5.5 节引入的术语说，即
\begin{align}
\mathrm{D}\eta(U)P(U)\geqslant a|P(U)|^2,\qquad U\in\mathscr{O},
\label{align:ch16-s06-e01}
\end{align}
其中 $a>0$。于是熵不等式 readily 给出 $|U|^2$ 的空间积分与 $|P(U)|^2$ 的时空积分的界。

由于我们需要 $|U|^2$ 在时空积分上的界，必须调用第 5.5 节引入的、流与源之间的协同——由 Kawashima 条件体现——
\begin{align}
\mathrm{D}P(0)R_i(0)\neq0,\qquad i=1,\ldots,n,
\label{align:ch16-s06-e02}
\end{align}
它保证：在平衡态 $0$ 附近线性化\eqref{align:ch16-s02-e01}所得的系统，不允许形如 $u(x-\lambda_i(0)t)R_i(0)$ 的解，即不表现无阻尼传播锋面。

我们证明：在\eqref{align:ch16-s06-e01}存在下，Kawashima 条件等价于\cref{theorem:ch16-1631}的假设\eqref{align:ch16-s03-e02}。由\eqref{align:ch16-s06-e01}，函数
\begin{align}
\theta(U)=\mathrm{D}\eta(U)P(U)-a|P(U)|^2
\label{align:ch16-s06-e03}
\end{align}
在 $U=0$ 处取最小值，从而 Hessian 矩阵
\begin{align}
\mathrm{D}^2\theta(0)=\mathrm{D}^2\eta(0)\mathrm{D}P(0)+\mathrm{D}P(0)^\top\mathrm{D}^2\eta(0)-2a\mathrm{D}P(0)^\top\mathrm{D}P(0)
\label{align:ch16-s06-e04}
\end{align}
半正定。从左以 $R_i^\top$ 乘原书(7.4.2)，得到
\begin{align}
R_i^\top\mathrm{D}^2\eta\mathrm{D}F=\lambda_iR_i^\top\mathrm{D}^2\eta,
\label{align:ch16-s06-e05}
\end{align}
这表明 $R_i^\top\mathrm{D}^2\eta$ 与 $L_i$ 共线，特别地
\begin{align}
R_i^\top\mathrm{D}^2\eta=(R_i^\top\mathrm{D}^2\eta R_i)L_i.
\label{align:ch16-s06-e06}
\end{align}

现在从左以 $R_i^\top(0)$、从右以 $R_i(0)$ 乘\eqref{align:ch16-s06-e04}。利用 $\mathrm{D}^2\theta(0)$ 半正定，连同\eqref{align:ch16-s06-e05}在 $U=0$ 处，我们断定
\begin{align}
[R_i^\top(0)\mathrm{D}^2\eta(0)R_i(0)]A_{ii}\geqslant a|\mathrm{D}P(0)R_i(0)|^2,
\label{align:ch16-s06-e07}
\end{align}
这验证了\eqref{align:ch16-s06-e02}蕴涵\eqref{align:ch16-s03-e02}。另一方面，从\eqref{align:ch16-s02-e02}易见 $A_{ii}\neq0$（$i=1,\ldots,n$）蕴涵 Kawashima 条件。

下一个任务是识别并暴露状态向量中直接受阻尼影响的分量。设 $\ker\mathrm{D}P(0)$ 的维数为 $k$，$1\leqslant k<n$。于是 $\mathrm{D}P(0)=S^{-1}\Gamma S$，其中 $S$ 是非奇异 $n\times n$ 矩阵，$\Gamma$ 是形如
\begin{align}
\Gamma=\begin{bmatrix}0&0\\0&C\end{bmatrix},
\label{align:ch16-s06-e08}
\end{align}
的 $n\times n$ 矩阵，$C$ 是非奇异 $\ell\times\ell$ 矩阵，$\ell=n-k$。从 $U$ 换到新状态向量 $\hat U=SU$，配以 $\mathbb{R}^k$ 与 $\mathbb{R}^\ell$ 上的投影 $V$ 与 $W$，把\eqref{align:ch16-s02-e01}约化为系统
\begin{align}
\begin{cases}
\partial_tV+\partial_xG(V,W)+X(V,W)=0,\\
\partial_tW+\partial_xH(V,W)+CW+Y(V,W)=0,
\end{cases}
\label{align:ch16-s06-e09}
\end{align}
其中 $X,Y$ 及其一阶导数在原点为零。于是我们可以不失一般性地假设系统\eqref{align:ch16-s02-e01}具有形式\eqref{align:ch16-s06-e09}。在下文中，我们将视方便使用\eqref{align:ch16-s02-e01}或\eqref{align:ch16-s06-e09}。

正如第 5.5 节已指出的，应用中遇到的、刻画松弛现象的系统，一般由守恒律与平衡律结合而成，因而假设具有 $X\equiv0$ 的\eqref{align:ch16-s06-e09}形式。
在那种情形下，易见熵不等式\eqref{align:ch16-s06-e01}要求 $Y(V,0)=0$，从而\eqref{align:ch16-s06-e09}可写成更紧凑的形式
\begin{align}
\begin{cases}
\partial_tV+\partial_xG(V,W)=0,\\
\partial_tW+\partial_xH(V,W)+C(V,W)W=0.
\end{cases}
\label{align:ch16-s06-e10}
\end{align}
\eqref{align:ch16-s06-e10}解的分量 $V$ 的空间均值是时间不变的，这大大简化了分析。因此，我们将在下面详细讨论形如\eqref{align:ch16-s06-e10}的系统这一特殊情形，在本节结尾简要回到一般情形\eqref{align:ch16-s06-e09}。

以熵 $\eta$ 作为 $(V,W)$ 的函数，Hessian 矩阵\eqref{align:ch16-s06-e04}的半正定成立当且仅当 $k\times\ell$ 矩阵 $\eta_{VW}(0,0)$ 为零，且 $\ell\times\ell$ 矩阵 $\eta_{WW}(0,0)C(0,0)$ 正定。此外，当对 $i=1,\ldots,n$，$R_i(0)$ 在 $\mathbb{R}^\ell$ 上的投影不为零时，Kawashima 条件\eqref{align:ch16-s06-e02}满足。

我们已经看到上述条件蕴涵\eqref{align:ch16-s03-e02}。着眼于\cref{theorem:ch16-1631}的其余假设，我们进而识别生成具有性质\eqref{align:ch16-s03-e01}之解的初值类 $\mathscr{C}_\omega$。这将通过下述定理实现。

\begin{theorem}\label{theorem:ch16-1661}
假设系统\eqref{align:ch16-s02-e01}具有形式\eqref{align:ch16-s06-e10}，源相对于熵 $\eta$ 耗散半正定，且 Kawashima 条件成立。设 $U=(V,W)$ 是一个可容 $BV$ 解，初值 $U_0=(V_0,W_0)$，定义在带形 $(-\infty,\infty)\times[0,T)$ 上，取值于以原点为心、半径 $\rho$ 小的球 $\mathscr{B}_\rho$ 中。假设 $U_0$ 随 $|x|\to\infty$ 衰减得足够快，使积分
\begin{align}
\int_{-\infty}^{\infty}(1+x^2)|U_0(x)|^2\,\mathrm{d}x=\sigma^2
\label{align:ch16-s06-e11}
\end{align}
有限。此外，设
\begin{align}
\int_{-\infty}^{\infty}V_0(x)\,\mathrm{d}x=0.
\label{align:ch16-s06-e12}
\end{align}
则存在 $\sigma_0>0$（与 $T$ 无关），使得当 $\sigma<\sigma_0$ 时，
\begin{align}
\int_{-\infty}^{\infty}|U(x,t)|\,\mathrm{d}x\leqslant b\sigma,\qquad 0\leqslant t<T,
\label{align:ch16-s06-e13}
\end{align}
其中 $b$ 与 $T$ 无关。此外，若 $T=\infty$，
\begin{align}
\int_{-\infty}^{\infty}|U(x,t)|\,\mathrm{d}x\to0,\qquad\text{当 }t\to\infty.
\label{align:ch16-s06-e14}
\end{align}
\end{theorem}

正是这个定理保证 $\mathscr{C}_\omega$ 包含所有满足\eqref{align:ch16-s06-e02}与\eqref{align:ch16-s06-e11}且 $\sigma\leqslant\omega/b$ 的初值 $U_0$。证明将作为下面要建立的先验估计的顶点。如同\cref{section:ch16-变差界}那样，为避免符号激增，我们通篇用 $c$ 表示一个一般正常数，它可能依赖于 $F,P$ 及其导数在原点某固定邻域（包含 $\mathscr{B}_\rho$）上的界，但与 $\rho,\sigma,T$ 无关。

一个有用的工具是“势”函数
\begin{align}
\Psi(x,t)=\int_{-\infty}^x V(y,t)\,\mathrm{d}y.
\label{align:ch16-s06-e15}
\end{align}
显然，$\Psi$ 是 Lipschitz 的，导数为
\begin{align}
\partial_x\Psi=V,\qquad\partial_t\Psi=-G(V,W).
\label{align:ch16-s06-e16}
\end{align}
假设\eqref{align:ch16-s06-e12}的作用是保证 $\Psi$ 的初值 $\Psi_0(x)$ 满足 $\Psi_0(\pm\infty)=0$。于是分部积分，
\begin{align}
\int_{-\infty}^{\infty}|\Psi_0(x)|^2\,\mathrm{d}x&=\int_{-\infty}^0\left|\int_{-\infty}^x V_0(y)\,\mathrm{d}y\right|^2\mathrm{d}x+\int_0^\infty\left|\int_\infty^x V_0(y)\,\mathrm{d}y\right|^2\mathrm{d}x\nonumber\\
&=-2\int_{-\infty}^0 xV_0^\top(x)\Psi_0(x)\,\mathrm{d}x-2\int_0^\infty xV_0^\top(x)\Psi_0(x)\,\mathrm{d}x\nonumber\\
&\leqslant2\int_{-\infty}^\infty x^2|V_0(y)|^2\,\mathrm{d}x+\frac12\int_{-\infty}^\infty|\Psi_0(x)|^2\,\mathrm{d}x,
\label{align:ch16-s06-e17}
\end{align}
从而由\eqref{align:ch16-s06-e11}，
\begin{align}
\int_{-\infty}^{\infty}|\Psi_0(x)|^2\,\mathrm{d}x\leqslant4\sigma^2.
\label{align:ch16-s06-e18}
\end{align}

下述引理首次提示了解的空间 $L^2$ 范数随时间趋于无穷而衰减。

\begin{lemma}\label{lemma:ch16-1662}
在\cref{theorem:ch16-1661}的假设下，存在与 $T$ 无关的常数 $\bar\omega$，使得若
\begin{align}
|\Psi(x,t)|<\bar\omega,\qquad-\infty<x<\infty,\ 0\leqslant t<T,
\label{align:ch16-s06-e19}
\end{align}
则
\begin{align}
\int_0^T\int_{-\infty}^{\infty}|U(x,t)|^2\,\mathrm{d}x\,\mathrm{d}t\leqslant c\sigma^2.
\label{align:ch16-s06-e020}
\end{align}
\end{lemma}

\begin{proof}
由\eqref{align:ch16-s06-e01}，把熵不等式\eqref{align:ch16-s05-e01}在 $(-\infty,\infty)\times[0,T)$ 上积分， readily 给出 $\int\int|W|^2\,\mathrm{d}x\,\mathrm{d}t$ 由 $c\sigma^2$ 界住。然而，要证明互补分量 $V$ 也是如此，则需要相当大的努力，因为它依赖于部分耗散源与流之间的协同——由 Kawashima 条件编码。

我们引入下述记号：
\begin{align}
B=G_V(0,0),\quad J=G_W(0,0),\quad E=H_V(0,0),\quad D=H_W(0,0),
\label{align:ch16-s06-e021}\\
K=\eta_{VV}(0,0),\quad M=\eta_{WW}(0,0),
\label{align:ch16-s06-e022}\\
Q=[\eta_{WW}(0,0)C(0,0)]^{-1}.
\label{align:ch16-s06-e023}
\end{align}
$k\times k$ 矩阵 $K$ 与 $\ell\times\ell$ 矩阵 $M$ 对称正定。此外，由于 $\mathrm{D}^2\eta(U)\mathrm{D}F(U)$ 对称且 $\eta_{VW}(0,0)=0$，
\begin{align}
(KB)^\top=KB,\quad(MD)^\top=MD,\quad(ME)^\top=KJ.
\label{align:ch16-s06-e024}
\end{align}
最后，$\ell\times\ell$ 矩阵 $Q$ 正定。

若 $N$ 是矩阵 $B$ 的任一特征向量，则 $EN\neq0$，因为 $EN=0$ 将蕴涵 $R=\begin{pmatrix}N\\0\end{pmatrix}$ 是 $DF(0)$ 的特征向量且 $\mathrm{D}P(0)R=0$，与 Kawashima 条件矛盾。于是可证明（见\cref{section:ch16-注记}所引文献）存在 $k\times k$ 矩阵 $\Omega$，使 $\Omega K$ 斜对称且 $\Omega KB$ 在 $ME$ 的核上为正。

我们现在定义下述函数：
\begin{align}
\Theta(V,W,\Psi)&=\Psi^\top K\Psi-2\Psi^\top KJQMW-\kappa\Psi^\top\Omega KV+\gamma\eta(V,W),
\label{align:ch16-s06-e025}\\
\Xi(V,W,\Psi)&=\Psi^\top KB\Psi-2\Psi^\top KJQMH(V,W)-\kappa\Psi^\top\Omega KG(V,W)+\gamma q(V,W),
\label{align:ch16-s06-e026}\\
\Pi(V,W,\Psi)&=2\Psi^\top KG(V,W)-2\Psi^\top KBV-2G^\top(V,W)KJQMW\nonumber\\
&\quad+2V^\top KJQMH(V,W)-2\Psi^\top KJQMC(V,W)W-\kappa G^\top(V,W)\Omega KV\nonumber\\
&\quad+\kappa V^\top\Omega KG(V,W)+\gamma\eta_W(V,W)C(V,W)W,
\label{align:ch16-s06-e027}
\end{align}
其中 $\kappa$ 与 $\gamma$ 是待在下面固定的正常数。

一个冗长但直接的计算，使用\eqref{align:ch16-s06-e10}、\eqref{align:ch16-s06-e16}与\eqref{align:ch16-s05-e01}，给出
\begin{align}
\partial_t\Theta(V,W,\Psi)+\partial_x\Xi(V,W,\Psi)+\Pi(V,W,\Psi)\leqslant0.
\label{align:ch16-s06-e028}
\end{align}

我们对 $\Pi(V,W,\Psi)$ 关于原点作（有限）Taylor 展开。利用对称关系\eqref{align:ch16-s06-e024}，并回忆 $|U|<\rho$、$|\Psi|<\bar\omega$，我们得到
\begin{align}
\Pi=V^\top\Lambda V+V^\top\Gamma W+W^\top\Delta W+O(\rho+\bar\omega)(|V|^2+|W|^2),
\label{align:ch16-s06-e029}
\end{align}
其中
\begin{align}
\Lambda=2(ME)^\top QME+2\kappa\Omega KB,
\label{align:ch16-s06-e30}\\
\Delta=-2J^\top KJQM+\gamma Q^{-1},
\label{align:ch16-s06-e31}\\
\Gamma=-2KBJQM+2KJQMD+2\kappa\Omega KJ.
\label{align:ch16-s06-e32}
\end{align}

关键观察是：\eqref{align:ch16-s06-e30}右端的第二项在 $ME$ 的核上为正，第一项在补空间上为正，从而当 $\kappa$ 充分小时，$\Lambda$ 正定。于是我们可以取 $\rho$ 与 $\bar\omega$ 充分小、$\gamma$ 充分大，使 $\Theta$ 与 $\Pi$ 在原点都正定。然后在 $(-\infty,\infty)\times[0,T)$ 上积分\eqref{align:ch16-s06-e028}，并使用\eqref{align:ch16-s06-e11}与\eqref{align:ch16-s06-e18}，得到\eqref{align:ch16-s06-e020}。证明完成。
\end{proof}

\begin{lemma}\label{lemma:ch16-1663}
在\cref{theorem:ch16-1661}的假设下，只要\eqref{align:ch16-s06-e19}成立，
\begin{align}
t\int_{-\infty}^{\infty}|U(x,t)|^2\,\mathrm{d}x\leqslant c\sigma^2,\qquad 0\leqslant t<T.
\label{align:ch16-s06-e33}
\end{align}
此外，若 $T=\infty$，
\begin{align}
t\int_{-\infty}^{\infty}|U(x,t)|^2\,\mathrm{d}x\to0,\qquad\text{当 }t\to\infty.
\label{align:ch16-s06-e34}
\end{align}
\end{lemma}

\begin{proof}
由\eqref{align:ch16-s05-e01}、\eqref{align:ch16-s05-e02}，
\begin{align}
\partial_t[t\eta(U)]+\partial_x[tq(U)]\leqslant\eta(U).
\label{align:ch16-s06-e35}
\end{align}
把上式在 $(-\infty,\infty)\times[0,t]$ 上积分（$t\in[0,T)$），并使用\eqref{align:ch16-s06-e020}，我们得到\eqref{align:ch16-s06-e33}。

现设 $T=\infty$，固定任意 $\varepsilon>0$。由\eqref{align:ch16-s06-e020}，存在 $\tau>0$ 使得
\begin{align}
\tau\int_{-\infty}^{\infty}\eta(U(x,\tau))\,\mathrm{d}x<\frac{\varepsilon}{2},
\label{align:ch16-s06-e36}\\
\int_\tau^\infty\int_{-\infty}^{\infty}\eta(U(x,t))\,\mathrm{d}x\,\mathrm{d}t<\frac{\varepsilon}{2}.
\label{align:ch16-s06-e37}
\end{align}
对任意 $t\in(\tau,\infty)$，在 $(-\infty,\infty)\times[\tau,t]$ 上积分\eqref{align:ch16-s06-e35}，得到
\begin{align}
t\int_{-\infty}^{\infty}\eta(U(x,t))\,\mathrm{d}x<\varepsilon,
\label{align:ch16-s06-e38}
\end{align}
由此得\eqref{align:ch16-s06-e34}。证明完成。
\end{proof}

\textbf{\cref{theorem:ch16-1661}的证明。}假设\eqref{align:ch16-s06-e19}成立。我们确定 $\lambda>0$，使得对包含 $\mathscr{B}_\rho$ 的原点某邻域中一切 $U$，
\begin{align}
|q(U)|\leqslant\lambda\eta(U)
\label{align:ch16-s06-e39}
\end{align}
成立。

由\eqref{align:ch16-s06-e33}与 Schwarz 不等式，
\begin{align}
\int_{-2\lambda t}^{2\lambda t}|U(x,t)|\,\mathrm{d}x\leqslant\left\{4\lambda t\int_{-2\lambda t}^{2\lambda t}|U(x,t)|^2\,\mathrm{d}x\right\}^{1/2}\leqslant c\sigma,
\label{align:ch16-s06-e40}
\end{align}
对任意 $t\in[0,T)$。此外，若 $T=\infty$，\eqref{align:ch16-s06-e40}与\eqref{align:ch16-s06-e34}蕴涵
\begin{align}
\int_{-2\lambda t}^{2\lambda t}|U(x,t)|\,\mathrm{d}x\to0,\qquad\text{当 }t\to\infty.
\label{align:ch16-s06-e41}
\end{align}

我们进而估计 $|U(x,t)|$ 在区间 $(2\lambda t,\infty)$ 上的积分。对 $k=1,2,\ldots$，把熵不等式\eqref{align:ch16-s05-e01}在顶点为 $(2^k\lambda,t)$、$(2^{k+1}\lambda,t)$、$((2^k-1)\lambda,0)$ 与 $((2^{k+1}+1)\lambda,0)$ 的梯形上积分。由\eqref{align:ch16-s06-e39}，这给出
\begin{align}
\int_{2^k\lambda t}^{2^{k+1}\lambda t}|U(x,t)|^2\,\mathrm{d}x\leqslant\beta\int_{(2^k-1)\lambda t}^{(2^{k+1}+1)\lambda t}|U_0(x)|^2\,\mathrm{d}x.
\label{align:ch16-s06-e42}
\end{align}
因此，
\begin{align}
\int_{2^k\lambda t}^{2^{k+1}\lambda t}(x^2+1)|U(x,t)|^2\,\mathrm{d}x&\leqslant\beta(4^{k+1}\lambda^2t^2+1)\int_{(2^k-1)\lambda t}^{(2^{k+1}+1)\lambda t}|U_0(x)|^2\,\mathrm{d}x\nonumber\\
&\leqslant16\beta\int_{(2^k-1)\lambda t}^{(2^{k+1}+1)\lambda t}(x^2+1)|U_0(x)|^2\,\mathrm{d}x.
\label{align:ch16-s06-e43}
\end{align}
对 $k=1,2,\ldots$ 求和上述不等式，我们推得
\begin{align}
\int_{2\lambda t}^{\infty}(x^2+1)|U(x,t)|^2\,\mathrm{d}x\leqslant32\beta\int_{\lambda t}^{\infty}(x^2+1)|U_0(x)|^2\,\mathrm{d}x\leqslant c\sigma^2.
\label{align:ch16-s06-e44}
\end{align}
最后，由 Schwarz 不等式，
\begin{align}
\int_{2\lambda t}^{\infty}|U(x,t)|\,\mathrm{d}x\leqslant\left[\int_{2\lambda t}^{\infty}(x^2+1)^{-1}\,\mathrm{d}x\right]^{1/2}\left[\int_{2\lambda t}^{\infty}(x^2+1)|U_0(x)|^2\,\mathrm{d}x\right]^{1/2}\leqslant c\sigma(\lambda t+1)^{-1/2}.
\label{align:ch16-s06-e45}
\end{align}
对 $|U(x,t)|$ 在 $(-\infty,-2\lambda t)$ 上的积分得到类似的界。这些界连同\eqref{align:ch16-s06-e40}与\eqref{align:ch16-s06-e41}，对某个 $b$ 建立了\eqref{align:ch16-s06-e13}以及\eqref{align:ch16-s06-e14}，但受制于\eqref{align:ch16-s06-e19}。为证明这一限制是多余的，我们首先注意：必要时增大 $b$ 的大小，\eqref{align:ch16-s06-e13}在 $t=0$ 处成立，独立于\eqref{align:ch16-s06-e19}。于是我们取 $\sigma_0=\bar\omega/b$ 并固定 $\sigma<\sigma_0$。由于 $|\Psi(x,t)|\leqslant\|V(\cdot,t)\|_{L^1}$，我们推知：当\eqref{align:ch16-s06-e13}对某个 $t\in[0,T)$ 满足时，由\eqref{align:ch16-s06-e16}，\eqref{align:ch16-s06-e19}必在 $(-\infty,\infty)\times[t+\varepsilon)$ 上成立。一个简单的延拓论证然后建立：对任意 $\sigma<\sigma_0$，\eqref{align:ch16-s06-e19}在 $(-\infty,\infty)\times[0,T)$ 上成立。证明完成。

现在我们可以合并\cref{theorem:ch16-1631}与\cref{theorem:ch16-1661}，得到下述存在性定理：

\begin{theorem}\label{theorem:ch16-1664}
考虑形如\eqref{align:ch16-s06-e10}的平衡律组\eqref{align:ch16-s02-e01}，其特征场真非线性或线性退化，配有凸熵 $\eta$。假设源相对于 $\eta$ 耗散半正定\eqref{align:ch16-s06-e01}，且 Kawashima 条件\eqref{align:ch16-s06-e02}成立。则存在正常数 $\delta_0,\sigma_0,c_0,c_1,\nu$ 与 $b$，使得在初值 $U_0=(V_0,W_0)$ 满足
\begin{align}
\int_{-\infty}^{\infty}(1+x^2)|U_0(x)|^2\,\mathrm{d}x=\sigma^2<\sigma_0^2,
\label{align:ch16-s06-e46}\\
TV_{(-\infty,\infty)}U_0(\cdot)=\delta<\delta_0,
\label{align:ch16-s06-e47}\\
\int_{-\infty}^{\infty}V_0(x)\,\mathrm{d}x=0
\label{align:ch16-s06-e48}
\end{align}
时，Cauchy 问题在 $(-\infty,\infty)\times[0,\infty)$ 上拥有一个可容 $BV$ 解 $U$，且
\begin{align}
\int_{-\infty}^{\infty}|U(x,t)|\,\mathrm{d}x\leqslant b\sigma,\qquad 0\leqslant t<\infty,
\label{align:ch16-s06-e49}\\
TV_{(-\infty,\infty)}U(\cdot,t)\leqslant c_0\sigma+c_1\delta e^{-\nu t},\qquad 0\leqslant t<\infty,
\label{align:ch16-s06-e50}\\
\int_{-\infty}^{\infty}|U(x,t)|\,\mathrm{d}x\to0,\qquad\text{当 }t\to\infty,
\label{align:ch16-s06-e51}\\
TV_{(-\infty,\infty)}U(\cdot,t)\to0,\qquad\text{当 }t\to\infty.
\label{align:ch16-s06-e52}
\end{align}
\end{theorem}

作为说明，我们再次考虑\cref{section:ch16-原点锥奇异熵l1稳定}中讨论的系统原书(5.5.52)。
为在本节的语境下处理它，我们需要一个凸熵 $\eta$，使源耗散半正定：
\begin{align}
\eta_v(u,v)[v-f(u)]\geqslant\frac{a}{\mu}|v-f(u)|^2,
\label{align:ch16-s06-e53}
\end{align}
其中 $a>0$。

遵循\cref{section:ch16-原点锥奇异熵l1稳定}中的程序，我们将构造一个具有上述规格的熵，把它作为 Riemann 不变量\eqref{align:ch16-s05-e07}（定义在矩形\eqref{align:ch16-s05-e11}上）的函数 $\eta(z,w)$，方法是在 Cauchy 数据\eqref{align:ch16-s05-e12}、\eqref{align:ch16-s05-e13}下求解方程\eqref{align:ch16-s05-e08}的 Cauchy 问题。然而，此处我们使用任意凸函数 $\varphi$，满足 $\varphi(0)=0$、$\varphi_z(0)=0$、$\varphi_{zz}(0)=1$。我们执行与\cref{section:ch16-原点锥奇异熵l1稳定}中类似的估计，得到
\begin{align}
\eta(z,w)=O(r^2),
\label{align:ch16-s06-e54}\\
\eta_z(z,w)=O(r),\quad\eta_w(z,w)=O(r),\quad\eta_{zw}(z,w)=O(r),
\label{align:ch16-s06-e55}\\
\eta_{zz}(z,w)=\frac{1}{1+g_z(z)}+O(r),
\label{align:ch16-s06-e56}\\
\eta_{ww}(z,w)=\frac{1}{g_z(g^{-1}(w))[1+g_z(g^{-1}(w))]}+O(r),
\label{align:ch16-s06-e57}
\end{align}
此处阶符号 $O$ 绑定于假设 $\varphi_{zz}(0)=1$。

于是，当 $r$ 小时，$\eta$ 满足凸性条件\eqref{align:ch16-s05-e09}。此外，由于 $\eta_v=\eta_w-\eta_z$、$\eta_{vv}=\eta_{zz}+\eta_{ww}-2\eta_{zw}$，我们推得 $\eta_v(u,f(u))=0$ 且 $\eta_{vv}(u,v)\geqslant a/\mu$，从而\eqref{align:ch16-s06-e53}确实满足。

给定 $(-\infty,\infty)$ 上的函数 $u_0$ 与 $v_0$，满足
\begin{align}
\int_{-\infty}^{\infty}(1+x^2)[u_0^2(x)+v_0^2(x)]\,\mathrm{d}x=\sigma^2,
\label{align:ch16-s06-e58}\\
TV_{(-\infty,\infty)}u_0(\cdot)+TV_{(-\infty,\infty)}v_0(\cdot)=\delta,
\label{align:ch16-s06-e59}\\
\int_{-\infty}^{\infty}u_0(x)\,\mathrm{d}x=0,
\label{align:ch16-s06-e60}
\end{align}
其中 $\sigma$ 与 $\delta$ 充分小，\cref{theorem:ch16-1664}蕴涵：系统原书(5.5.52)以 $(u_0,v_0)$ 为初值的 Cauchy 问题在大范围拥有一个可容 $BV$ 解 $(u,v)$，且
\begin{align}
\int_{-\infty}^{\infty}[|u(x,t)|+|v(x,t)|]\,\mathrm{d}x\to0,\qquad\text{当 }t\to\infty,
\label{align:ch16-s06-e61}\\
TV_{(-\infty,\infty)}u(\cdot,t)+TV_{(-\infty,\infty)}v(\cdot,t)\to0,\qquad\text{当 }t\to\infty.
\label{align:ch16-s06-e62}
\end{align}

把上述与\cref{section:ch16-原点锥奇异熵l1稳定}中对系统原书(5.5.52)的处理相比较，人们注意到假设\eqref{align:ch16-s06-e58}比\eqref{align:ch16-s05-e20}更严格，且此处施加的条件\eqref{align:ch16-s06-e60}在那里并不需要。
另一方面，与\cref{section:ch16-原点锥奇异熵l1稳定}的途径相反，此处的处理还给出解的长时间行为。

事实上，对更一般的系统类\eqref{align:ch16-s06-e09}（带非零 $X$）的 Cauchy 问题、甚至对没有条件\eqref{align:ch16-s06-e12}的较窄类\eqref{align:ch16-s06-e10}，其 $BV$ 解的存在性与长时间行为问题仍未解决。显然，$U$——或至少其 $V$ 分量——一般不会在 $L^1(-\infty,\infty)$ 中随 $t\to\infty$ 衰减到零。下述特例结果首次提示了在这种情形下应期待什么。

我们考虑简单系统
\begin{align}
\begin{cases}
\partial_tu-\partial_xv+\alpha v^2=0,\\
\partial_tv+\partial_xp(u)+v=0,
\end{cases}
\label{align:ch16-s06-e63}
\end{align}
的 Cauchy 问题，作为类\eqref{align:ch16-s06-e09}中系统的模型，其中 $p'(u)<0$，$\alpha$ 是常数，可正、可负或为零。下述与\cref{theorem:ch16-1664}类似的定理成立：

\begin{theorem}\label{theorem:ch16-1665}
设 $(u_0,v_0)$ 是 $(-\infty,\infty)$ 上给定的函数，满足\eqref{align:ch16-s06-e58}与\eqref{align:ch16-s06-e59}，其中 $\sigma$ 与 $\delta$ 充分小。则系统\eqref{align:ch16-s06-e63}以初值 $(u_0,v_0)$ 在 $(-\infty,\infty)\times[0,\infty)$ 上存在一个可容 $BV$ 解 $(u,v)$。此外，
\begin{align}
\int_{-\infty}^{\infty}[|u(x,t)-\theta(x,t)|+|v(x,t)|]\,\mathrm{d}x\leqslant b\sigma(t+1)^{-1/4},\qquad 0\leqslant t<\infty,
\label{align:ch16-s06-e64}\\
TV_{(-\infty,\infty)}u(\cdot,t)+TV_{(-\infty,\infty)}v(\cdot,t)\leqslant c_0\sigma(t+1)^{-1/4}+c_1\delta e^{-\nu t},\qquad 0\leqslant t<\infty,
\label{align:ch16-s06-e65}
\end{align}
其中 $c_0,c_1,\nu$ 与 $b$ 是正常数，与初值无关，而 $\theta$ 是热方程
\begin{align}
\theta(x,t)=M(4\pi t)^{-1/2}\exp\left(-\frac{x^2}{4t}\right)
\label{align:ch16-s06-e66}
\end{align}
的解，其中 $M$ 是依赖于 $(u_0,v_0)$ 的某个常数。
\end{theorem}

\textbf{证明梗概。}为简单起见，我们只讨论 $\alpha=0$ 的特殊情形，此时\eqref{align:ch16-s06-e63}仍具有\eqref{align:ch16-s06-e10}的形式。然而，我们不再施加\eqref{align:ch16-s06-e60}，而是假设
\begin{align}
\int_{-\infty}^{\infty}u_0(x)\,\mathrm{d}x=M.
\label{align:ch16-s06-e67}
\end{align}
我们仅仅勾勒证明。细节连同一般情形 $\alpha\neq0$ 的处理，见\cref{section:ch16-注记}所引文献。

目标是证明：对任意 $0<t<\infty$，
\begin{align}
\int_{-\infty}^{\infty}[|u(x,t)-\hat u(x,t)|+|v(x,t)-\hat v(x,t)|]\,\mathrm{d}x\leqslant c\sigma(t+1)^{-1/4},
\label{align:ch16-s06-e68}
\end{align}
其中 $(\hat u,\hat v)$ 是系统
\begin{align}
\begin{cases}
\partial_t\hat u-\partial_x\hat v=0,\\
\hat v=-\partial_xp(\hat u),
\end{cases}
\label{align:ch16-s06-e69}
\end{align}
以与 $(u,v)$ 相同的初值 $(u_0,v_0)$ 的解。于是 $\hat u$ 满足多孔介质方程
\begin{align}
\partial_t\hat u+\partial_x^2p(\hat u)=0.
\label{align:ch16-s06-e70}
\end{align}
借助冗长但直接的分析（涉及初等“能量”估计），可以建立\eqref{align:ch16-s06-e70}解 $\hat u$ 并从而 $\hat v$ 的界，包括下列：
\begin{align}
\int_0^\infty\int_{-\infty}^{\infty}\hat v^2(x,t)\,\mathrm{d}x\,\mathrm{d}t\leqslant c\sigma^2,
\label{align:ch16-s06-e71}\\
\int_0^\infty\int_{-\infty}^{\infty}(x^2+t+1)[\partial_t\hat u(x,t)]^2\,\mathrm{d}x\,\mathrm{d}t\leqslant c\sigma^2,
\label{align:ch16-s06-e72}\\
\int_0^\infty\int_{-\infty}^{\infty}(x^2+t+1)^2[\partial_t\hat v(x,t)]^2\,\mathrm{d}x\,\mathrm{d}t\leqslant c\sigma^2,
\label{align:ch16-s06-e73}\\
\int_{-\infty}^{\infty}|\hat u(x,t)-\theta(x,t)|\,\mathrm{d}x\leqslant c\sigma(t+1)^{-1/4},
\label{align:ch16-s06-e74}\\
\int_{-\infty}^{\infty}|\hat v(x,t)|\,\mathrm{d}x\leqslant c\sigma(t+1)^{-1/2},
\label{align:ch16-s06-e75}
\end{align}
其中 $\theta$ 由\eqref{align:ch16-s06-e66}定义。于是证明\eqref{align:ch16-s06-e68}将建立断言\eqref{align:ch16-s06-e64}。变差中所断言的衰减\eqref{align:ch16-s06-e65}然后由\cref{theorem:ch16-1631}得到。

着眼于验证\eqref{align:ch16-s06-e68}，我们令 $w=u-\hat u$、$z=v-\hat v$，注意到 $(w,z)$ 满足系统
\begin{align}
\begin{cases}
\partial_tw-\partial_xz=0,\\
\partial_tz+\partial_x\hat p(w,\hat u)+z+\partial_t\hat v=0,
\end{cases}
\label{align:ch16-s06-e76}
\end{align}
带零初值。在\eqref{align:ch16-s06-e76}中，$\hat p$ 表示由
\begin{align}
\hat p(w,\hat u)=p(w+\hat u)-p(\hat u)
\label{align:ch16-s06-e77}
\end{align}
定义的“相对压力”。

\eqref{align:ch16-s06-e76}解的可容性编码在“相对熵”不等式中
\begin{align}
\partial_t\left[\hat\psi(w,\hat u)+\frac12z^2\right]+\partial_x[\hat p(w,\hat u)z]+z^2\leqslant-[\hat p(w,\hat u)-p'(\hat u)w]\partial_t\hat u-z\partial_t\hat v,
\label{align:ch16-s06-e78}
\end{align}
其中 $\hat\psi$ 表示由
\begin{align}
\hat\psi(w,\hat u)=-\int_{\hat u}^{w+\hat u}p(\xi)\,\mathrm{d}\xi+p(\hat u)w
\label{align:ch16-s06-e79}
\end{align}
定义的“相对内能”。

$(w,z)$ 的先验界将通过把\eqref{align:ch16-s06-e78}与平衡律
\begin{align}
\partial_t\left[\frac12\Phi^2+(z+\hat v)\Phi\right]+\partial_x[\hat p(w,\hat u)\Phi]-\hat p(w,\hat u)w=z^2+z\hat v,
\label{align:ch16-s06-e80}
\end{align}
相结合而导出，其中 $\Phi$ 是“势函数”
\begin{align}
\Phi(x,t)=\int_{-\infty}^x w(y,t)\,\mathrm{d}y.
\label{align:ch16-s06-e81}
\end{align}
\eqref{align:ch16-s06-e78}与\eqref{align:ch16-s06-e80}中出现的具“好符号”的项包括 $\hat\psi(w,\hat u)$、$z^2$、$\Phi^2$ 与 $-\hat p(w,\hat u)w$。易见，当 $\sigma$ 充分小时，借助\eqref{align:ch16-s06-e71}、\eqref{align:ch16-s06-e72}与\eqref{align:ch16-s06-e73}，不定号的项可与具好符号的项相平衡，从而给出界
\begin{align}
\int_{-\infty}^{\infty}[w^2(x,t)+z^2(x,t)]\,\mathrm{d}x\leqslant c\sigma^2,
\label{align:ch16-s06-e82}\\
\int_0^\infty\int_{-\infty}^{\infty}[w^2(x,t)+z^2(x,t)]\,\mathrm{d}x\,\mathrm{d}t\leqslant c\sigma^2.
\label{align:ch16-s06-e83}
\end{align}

为得到下一轮估计，我们先以 $t$、再以 $x^2$ 乘\eqref{align:ch16-s06-e78}，这给出
\begin{align}
\partial_t\left[t\left[\hat\psi(w,\hat u)+\frac12z^2\right]\right]+\partial_x[t\hat p(w,\hat u)z]+tz^2\nonumber\\
\qquad\leqslant\hat\psi(w,\hat u)+\frac12z^2-t[\hat p(w,\hat u)-p'(\hat u)w]\partial_t\hat u-tz\partial_t\hat v,
\label{align:ch16-s06-e84}\\
\partial_t\left[x^2\left[\hat\psi(w,\hat u)+\frac12z^2\right]\right]+\partial_x[x^2\hat p(w,\hat u)z]+x^2z^2\nonumber\\
\qquad\leqslant2x\hat p(w,\hat u)z-x^2[\hat p(w,\hat u)-p'(\hat u)w]\partial_t\hat u-x^2z\partial_t\hat v.
\label{align:ch16-s06-e85}
\end{align}
借助\eqref{align:ch16-s06-e82}与\eqref{align:ch16-s06-e83}，连同\eqref{align:ch16-s06-e71}、\eqref{align:ch16-s06-e72}与\eqref{align:ch16-s06-e73}，可以把\eqref{align:ch16-s06-e84}与\eqref{align:ch16-s06-e85}中的不定号项与具好符号的项 $t\hat\psi$、$x^2\hat\psi$、$tz^2$ 与 $x^2z^2$ 相平衡，从而得到估计
\begin{align}
\int_{-\infty}^{\infty}(x^2+t+1)[w^2(x,t)+z^2(x,t)]\,\mathrm{d}x\leqslant c\sigma^2.
\label{align:ch16-s06-e86}
\end{align}
最后，对\eqref{align:ch16-s06-e86}应用 Schwarz 不等式，得到
\begin{align}
\left[\int_{-\infty}^{\infty}[|w(x,t)|+|z(x,t)|]\,\mathrm{d}x\right]^2\nonumber\\
\qquad\leqslant2\int_{-\infty}^{\infty}(x^2+t+1)^{-1}\,\mathrm{d}x\int_{-\infty}^{\infty}(x^2+t+1)[w^2(x,t)+z^2(x,t)]\,\mathrm{d}x\nonumber\\
\qquad\leqslant c\sigma^2(t+1)^{-1/2},
\label{align:ch16-s06-e87}
\end{align}
由此我们得到\eqref{align:ch16-s06-e68}。证明梗概完成。

在本章所考虑的平衡律组类的 Cauchy 问题解的长时间行为中，当 $|x|\to\infty$ 初值衰减时，源的阻尼作用被波的色散所缓和。相反，当波被束缚——如同空间周期解的情形——源的耗散效应变成绝对主导。下述命题提供了一个例证，其证明见\cref{section:ch16-注记}所引文献。

\begin{theorem}\label{theorem:ch16-1666}
考虑形如\eqref{align:ch16-s06-e10}的平衡律组\eqref{align:ch16-s02-e01}，其特征场真非线性或线性退化，配有凸熵 $\eta$。假设源相对于 $\eta$ 耗散半正定\eqref{align:ch16-s06-e01}，且 Kawashima 条件\eqref{align:ch16-s06-e02}成立。则存在正常数 $\delta_0,c_0,c_1$ 与 $\nu$，使得当初值 $U_0$ 满足
\begin{align}
U_0(x+2)=U_0(x),\qquad-\infty<x<\infty,
\label{align:ch16-s06-e88}\\
|U_0(x)|<\delta<\delta_0,\qquad-\infty<x<\infty,
\label{align:ch16-s06-e89}\\
TV_{[-1,1]}U_0(\cdot)=\delta<\delta_0
\label{align:ch16-s06-e90}
\end{align}
时，Cauchy 问题在 $(-\infty,\infty)\times[0,\infty)$ 上拥有一个可容 $BV$ 解 $U$，它是 $2$-周期的，并满足
\begin{align}
|U(x,t)|\leqslant c_0\delta e^{-\nu t},\qquad-1\leqslant x<1,\ 0\leqslant t<\infty,
\label{align:ch16-s06-e91}\\
TV_{[-1,1]}U(\cdot,t)\leqslant c_1\delta e^{-\nu t},\qquad 0\leqslant t<\infty.
\label{align:ch16-s06-e92}
\end{align}
\end{theorem}

特别地，在严格不等式形式的次特征条件原书(5.5.54)下，系统原书(5.5.52)满足\cref{theorem:ch16-1666}的条件，从而拥有随时间趋于无穷而指数衰减的空间周期 $BV$ 解。

真非线性两守恒律组的空间周期 $BV$ 解的存在性已在第十二章建立；然而，它是否推广到规模更大的守恒律组是令人怀疑的。
我们在此看到，在耗散源存在下，这些困难消失了。

\end{document}
